%=====================================================================
% PART I OPENER
%=====================================================================
\thispagestyle{plain}
\structuralfigures{P1.}

\begin{center}
\begin{tcolorbox}[enhanced, width=0.92\textwidth, colback=primary!5,
  colframe=primary, boxrule=0pt, leftrule=5pt, arc=3pt, top=8pt, bottom=8pt]
\large\itshape\color{primary!80!black}
Every idea in this book is a special case of one move: put a layer underneath
something, and make what sits above believe nothing has changed. Part~I says
what that layer is, names the ways it can be built --- and then proves the
uncomfortable fact that for twenty years the most widely sold processor in the
world could not host one correctly.
\end{tcolorbox}
\end{center}

\vspace{6pt}

\dsfigh{%
\begin{tikzpicture}[
  cbox/.style={rounded corners=3pt, draw=primary, line width=1pt, fill=primary!7,
               text=primary!80!black, font=\small\bfseries, align=center,
               minimum width=3.6cm, minimum height=1.0cm},
  hbox/.style={rounded corners=3pt, draw=primary, line width=1.8pt, fill=primary!16,
               text=primary!90!black, font=\small\bfseries, align=center,
               minimum width=3.6cm, minimum height=1.0cm},
  arr/.style={-{Stealth[length=2.4mm]}, primary!60, line width=1pt}
]
  \node[cbox] (c1) at (0,0)   {Ch.\ 1\\What Virtualization Is};
  \node[cbox] (c2) at (4.6,0) {Ch.\ 2\\Architectures};
  \node[hbox] (c3) at (9.2,0) {Ch.\ 3\\The Virtual\\Machine Monitor};
  \draw[arr] (c1) -- (c2);
  \draw[arr] (c2) -- (c3);
  \node[font=\scriptsize\itshape, text=primary, align=center, fill=white,
        inner sep=2pt] at (2.3,-1.0) {how it can be built};
  \node[font=\scriptsize\itshape, text=primary, align=center, fill=white,
        inner sep=2pt] at (6.9,-1.0) {when it is correct};
  \node[font=\scriptsize\itshape, text=accent, align=center] at (9.2,1.05)
       {the theory the whole of Part II answers};
\end{tikzpicture}}%
{Reading path through Part~I. Chapter~3 is the hinge: it states the conditions a
correct hypervisor must meet, and every hardware feature in Part~II exists
because x86 did not meet them.}{part1map}

\newpage
